\documentclass[aspectratio=169,10pt]{beamer}
\usetheme[numbering=fraction,progressbar=frametitle]{metropolis}
\usepackage{fontspec}
\setsansfont{Latin Modern Sans}
\usepackage{amsmath,amssymb,mathtools,graphicx,booktabs}
\usepackage{polyglossia}
\setdefaultlanguage{french}
\setbeamertemplate{navigation symbols}{}
\setbeamerfont{frametitle}{size=\large}
\setbeamerfont{normal text}{size=\normalsize}
\setlength{\parskip}{.45em}
\hypersetup{pdfauthor={},pdfsubject={Lecture slides},pdfkeywords={},pdftitle={Consolidation et transition vers la Première}}
\providecommand{\N}{\mathbb{N}}
\providecommand{\Z}{\mathbb{Z}}
\providecommand{\D}{\mathbb{D}}
\providecommand{\Q}{\mathbb{Q}}
\providecommand{\R}{\mathbb{R}}
\providecommand{\C}{\mathbb{C}}
\providecommand{\vect}[1]{\overrightarrow{#1}}
\providecommand{\vi}{\vec{\imath}}
\providecommand{\vj}{\vec{\jmath}}
\providecommand{\vk}{\vec{k}}
\providecommand{\norm}[1]{\left\lVert #1\right\rVert}
\providecommand{\vc}[2]{\begin{pmatrix}#1\\#2\end{pmatrix}}
\providecommand{\deter}[2]{\det\left(#1,#2\right)}
\providecommand{\abs}[1]{\left\lvert #1\right\rvert}
\providecommand{\dd}{\,\mathrm{d}}
\providecommand{\e}{\mathrm{e}}
\providecommand{\Card}{\operatorname{Card}}
\providecommand{\E}{\operatorname{E}}
\providecommand{\V}{\operatorname{V}}
\providecommand{\degre}{^{\circ}}
\providecommand{\un}[1]{\,\mathrm{#1}}
\providecommand{\sci}[2]{#1\times10^{#2}}
\providecommand{\pgcd}{\operatorname{PGCD}}
\providecommand{\ppcm}{\operatorname{PPCM}}
\title{Consolidation et transition vers la Première}
\subtitle{Cours et exemples}
\date{}
\author{}
\begin{document}
\begin{frame}[plain]
\vfill
{\usebeamerfont{title}\inserttitle\par}
\medskip
{\usebeamerfont{subtitle}\insertsubtitle\par}
\medskip
{\usebeamercolor[fg]{progress bar}\rule{\textwidth}{.5pt}}
\vfill
\end{frame}
\begin{frame}{Bilan de calcul, fonctions et géométrie}
\begin{enumerate}
\setcounter{enumi}{0}
\item Rendre \(168/252\) irréductible ; calculer \(2^{-3}+1/4\).

\item Factoriser \(x^2-6x+9\), puis résoudre \((x-3)^2=4\).

\item Résoudre \(\dfrac{x-2}{x+1}\geq0\).

\item Pour \(f(x)=4-(x-1)^2\) sur \([-2;3]\), donner le maximum et les abscisses où \(f(x)=0\).

\item Avec \(A(0;1)\), \(B(4;3)\), \(C(6;4)\), démontrer l'alignement et donner une équation de la droite.

\end{enumerate}
\end{frame}
\begin{frame}{Correction : Bilan de calcul, fonctions et géométrie}
1. \(168/252=2/3\), et \(2^{-3}+1/4=1/8+2/8=3/8\). 2. \((x-3)^2\) ; \(x-3=\pm2\), donc \(S=\{1;5\}\). 3. Interdit \(-1\), zéro \(2\) ; signes \(+,-,+\), donc \(S=]-\infty;-1[\cup[2;+\infty[\). 4. Le carré est positif ou nul, donc \(f(x)\leq4\), égalité en \(x=1\). Les zéros sont \(x=-1,3\), tous deux dans le domaine. 5. \(\vect{AB}=(4;2)\), \(\vect{AC}=(6;3)=\dfrac32\vect{AB}\). Une équation est \(x-2y+2=0\).
\end{frame}
\begin{frame}{Bilan de données, probabilités et modélisation}
\begin{enumerate}
\setcounter{enumi}{0}\item Calculer moyenne et médiane de \(2;4;4;6;9\). Déterminer \(Q_1,Q_3\).

\item Un prix baisse de \(15\,\%\), puis augmente de \(20\,\%\). Donner le taux global et le taux permettant ensuite de revenir au prix initial.

\item Dans une population fictive de \(100\) personnes, \(40\) font du sport (\(S\)), \(30\) viennent à vélo (\(V\)), \(18\) font les deux. Construire le tableau croisé et calculer \(P_S(V)\), \(P_V(S)\) pour un tirage équiprobable.

\end{enumerate}
\end{frame}
\begin{frame}{Bilan de données, probabilités et modélisation (2)}
\begin{enumerate}
\setcounter{enumi}{3}\item Pour un rectangle de périmètre \(24\) m, exprimer l'aire en fonction d'une dimension \(x\) et démontrer son maximum.

\end{enumerate}
\end{frame}
\begin{frame}{Correction : Bilan de données, probabilités et modélisation}
1. Moyenne \(25/5=5\), médiane \(4\). Rangs des quartiles \(2\) et \(4\), donc \(Q_1=4\), \(Q_3=6\).

2. Coefficient global \(0{,}85\times1{,}20=1{,}02\) : hausse de \(2\,\%\). Taux réciproque \(1/1{,}02-1=-1/51\approx-1{,}96\,\%\).

3. Tableau : \(S\cap V=18\), \(S\cap\overline V=22\), \(\overline S\cap V=12\), \(\overline S\cap\overline V=48\). \(P_S(V)=18/40=45\,\%\), \(P_V(S)=18/30=60\,\%\).

4. \(A(x)=x(12-x)=36-(x-6)^2\) sur \(]0;12[\). Le maximum vaut \(36\) m\(^2\), atteint pour le carré de côté \(6\) m.
\end{frame}
\begin{frame}{Choisir une stratégie et préparer la suite}
Pour une fonction donnée \(f\), expliquer le programme suivant. Que renvoie-t-il avec \(f(x)=x^2\) ?
\end{frame}
\begin{frame}{Correction : Choisir une stratégie et préparer la suite}
Il renvoie une liste de couples abscisse-image pour les entiers de \(-3\) à \(3\) inclus : \((-3;9),(-2;4),(-1;1),(0;0),(1;1),(2;4),(3;9)\). La borne \(4\) de \texttt{range} est exclue. Le tableau ne décrit pas les images de tous les réels entre \(-3\) et \(3\).
\end{frame}
\begin{frame}[fragile]{Python}
\small
\begin{verbatim}
def tableau(f):
    valeurs = []
    for x in range(-3, 4):
        valeurs.append((x, f(x)))
    return valeurs
\end{verbatim}
\end{frame}
\begin{frame}{Choisir une stratégie et préparer la suite}
\textbf{Chercher} : essayer, organiser des pistes. \textbf{Modéliser} : choisir variables, domaine et hypothèses. \textbf{Représenter} : passer du texte au tableau, à la figure ou au graphique. \textbf{Raisonner} : justifier et distinguer conjecture et preuve. \textbf{Calculer} : choisir une forme et contrôler les opérations. \textbf{Communiquer} : expliquer clairement avec des notations et unités adaptées.
\end{frame}
\begin{frame}{Choisir une stratégie et préparer la suite}
Choisir un travail réussi et un travail corrigé. Pour chacun, identifier les compétences mobilisées, expliquer une méthode et formuler un objectif personnel pour l'année suivante. Refaire ensuite un exercice correspondant à cet objectif sans regarder la correction.
\end{frame}
\begin{frame}{Correction : Choisir une stratégie et préparer la suite}
Il n'y a pas de réponse numérique unique. Une production satisfaisante identifie précisément le problème, la méthode choisie, une justification et un contrôle. L'objectif doit être observable, par exemple : « préciser les valeurs interdites avant de résoudre une équation quotient », plutôt que « devenir meilleur en calcul ».
\end{frame}
\end{document}
